\newpage
\topic{7 continued}{Prepare a steady iteration: under-relaxation}
A coupling \emph{iteration} updates an estimate of a steady solution. The
counter $k$ labels those updates. It does not measure elapsed physical time.
The estimates entering the update are $q^k$ and $p^k$.

\textbf{Velocity under-relaxation.} Convection coefficients depend on velocity,
so replacing an estimate by a very different predictor can destabilize the
next coefficient update. An under-relaxation factor $\alpha_u$, with
$0<\alpha_u\leq1$, controls this change. Subscript $u$ refers to velocity;
the same factor is used for both components here.

Let $a_f^0$ be the center coefficient before relaxation. The superscript $0$
identifies the original coefficient of this assembly, not an initial time.
The relaxed row uses
\[
a_f^{\rm rel}=\frac{a_f^0}{\alpha_u},\qquad
b_f^{\rm rel}=b_f+
\left(\frac{1}{\alpha_u}-1\right)a_f^0q_f^k.
\]
The label $\mathrm{rel}$ means relaxed. Neighbor coefficients are unchanged.
The larger diagonal is accompanied by a known source containing the old
velocity estimate. At a converged fixed point the new and old velocities
coincide, so these extra diagonal and source terms cancel. The intended
steady momentum equation is preserved.

For the preceding single-row example, let $a_e^0=0.64$, the old estimate be
$u_e^k=0.02$, and choose $\alpha_u=0.7$. Holding its neighbor and pressure
contributions fixed gives
\begin{align*}
a_e^{\rm rel}&=0.64/0.7\simeq0.9142857,\\
\text{added source}&=(1/0.7-1)(0.64)(0.02)\simeq0.005485714,\\
u_e^*&\simeq(0.0152+0.008+0.005485714)/0.9142857=0.031375.
\end{align*}
For this frozen-row illustration, this is also
$0.7(0.03625)+0.3(0.02)$. In the coupled solver, relaxation is applied
to the assembled momentum equations as written above. The pressure correction
uses the relaxed reciprocal diagonal $d_f=1/a_f^{\rm rel}$.

\textbf{Pressure under-relaxation.} The separate factor $\alpha_p$, with
$0<\alpha_p\leq1$, controls the pressure stored for the next predictor:
$p^{k+1}=p^k+\alpha_p p'$. The full velocity correction is retained to satisfy
continuity, as the two-cell example demonstrated. The factors $\alpha_u$ and
$\alpha_p$ are numerical controls, not fluid properties.
\newpage
\topic{7 continued}{Decide whether both equations have been satisfied}
A \emph{residual} is the left side of an equation minus its right side,
evaluated using the current estimates. Zero residual means that discrete
equation is satisfied. The cell continuity residual is already
$r_{c,P}=(Dq)_P$.

For the momentum equations, collect their residuals into a column $r_m$:
\[
r_m=A(q)q-b+Gp.
\]
Here $A(q)$ means that convection coefficients are rebuilt from the current
velocity before evaluating this test. Use the original, unrelaxed steady
equations. This checks the physical discretization being solved rather than
only the temporary relaxed predictor equations.

To report one number for a residual column $r$, use the infinity norm
\[
\|r\|_\infty=\max_i|r_i|.
\]
Index $i$ runs through that column's entries; the operation selects the
largest absolute entry. Thus the two required checks are
\[
\|r_m\|_\infty<\varepsilon_m,\qquad
\|Dq\|_\infty<\varepsilon_c.
\]
The positive numbers $\varepsilon_m$ and $\varepsilon_c$ are the selected
momentum and continuity tolerances. They are separate because the residuals
test different equations. The retained cavity runs use $10^{-6}$ and
$10^{-9}$ respectively.

A small change between successive velocities is insufficient: strong
under-relaxation can produce a small change even when the momentum balance
is still inaccurate. Small equation residuals are also distinct from small
spatial discretization error. Accuracy still requires the reference-profile
and mesh-refinement comparisons later in the lesson.

In a steady coupling iteration, both equation checks are evaluated after
the velocity and pressure updates, using the latest corrected fields.
